\chapter{Dise\~no e implementaci\'on del artefacto Axioma-H Engine}
\label{cap:diseno_artefacto}

\section{Arquitectura global del sistema}
El artefacto computacional \emph{Axioma-H Engine} se concibi\'o bajo el patr\'on de separaci\'on estricta entre el n\'ucleo computacional y la capa de presentaci\'on. La l\'ogica f\'isica, el motor de reglas y los algoritmos de an\'alisis operan en una API desacoplada desarrollada con FastAPI \parencite{fastapi}, mientras que el panel interactivo de usuario se implement\'o en Streamlit \parencite{streamlit} consumiendo la API v\'ia HTTP. La \cref{fig:arquitectura} ilustra el flujo de informaci\'on y la estructura modular.

\begin{figure}[htbp]
\centering
\begin{tikzpicture}[
  node distance=0.9cm and 1.1cm,
  bloque/.style={
    rectangle,
    rounded corners=3pt,
    draw=azulUdeA,
    thick,
    fill=grisClaro,
    text width=3.1cm,
    align=center,
    minimum height=1.3cm,
    font=\small
  },
  subbloque/.style={
    rectangle,
    rounded corners=2pt,
    draw=grisOscuro,
    fill=white,
    text width=2.8cm,
    align=center,
    minimum height=0.9cm,
    font=\footnotesize
  },
  flecha/.style={
    -{Stealth[length=2.5mm]},
    thick,
    draw=azulUdeA
  }
]

\node[bloque] (sensores) {\textbf{Telemetr\'ia y configuraci\'on}\\Sensores simulados ($T$, $\text{HR}$, $Z$, reserva) y par\'ametros};

\node[bloque, right=1.3cm of sensores, text width=3.6cm, minimum height=3.8cm, fill=white] (api) {};
\node[anchor=north, font=\small\bfseries, text=azulUdeA] at ([yshift=-0.15cm]api.north) {API REST (FastAPI :8000)};

\node[subbloque] (reglas) at ([yshift=0.7cm]api.center) {\textbf{Motor de reglas}\\(\texttt{motor\_reglas.py})};
\node[subbloque] (fisico) at ([yshift=-0.4cm]api.center) {\textbf{Modelo f\'isico y simulador}\\(\texttt{modelo\_fisico.py}, \texttt{simulador.py})};
\node[subbloque] (analisis) at ([yshift=-1.5cm]api.center) {\textbf{An\'alisis e impacto}\\(\texttt{analisis.py}, \texttt{impacto.py})};

\node[bloque, right=1.3cm of api] (dashboard) {\textbf{Dashboard (Streamlit :8501)}\\5 vistas interactivas\\2 niveles de lectura\\(\texttt{frontend/})};

\node[subbloque, below=0.8cm of api, text width=3.6cm] (datos) {\textbf{Datos abiertos y contrato}\\CSV, JSON, Data Package y OpenAPI 3.1 (\texttt{docs/})};

\draw[flecha] (sensores.east) -- node[above, font=\scriptsize] {JSON / HTTP} (api.west);
\draw[flecha] ([yshift=0.15cm]api.east) -- node[above, font=\scriptsize] {Respuestas JSON} ([yshift=0.15cm]dashboard.west);
\draw[flecha] ([yshift=-0.15cm]dashboard.west) -- node[below, font=\scriptsize] {Solicitudes HTTP} ([yshift=-0.15cm]api.east);
\draw[flecha] (api.south) -- node[right, font=\scriptsize] {Generaci\'on verificable} (datos.north);

\end{tikzpicture}


\caption{Arquitectura modular de Axioma-H Engine: desacoplamiento entre el backend anal\'itico en FastAPI y el frontend interactivo en Streamlit.}
\label{fig:arquitectura}
\end{figure}

\section{Motor de control reactivo basado en reglas deterministas}
El n\'ucleo de decisi\'on simb\'olico eval\'ua el vector de telemetr\'ia $\mathbf{t} = (T, \text{HR}, Z, E_{\text{bat}}, A_{\text{man}})$ seg\'un una jerarqu\'ia estricta de prioridades. La primera regla que se satisface determina la acci\'on de control ejecutada:
\begin{enumerate}[label=\textbf{Regla \arabic*:},leftmargin=*]
  \item \textbf{\texttt{REGLA\_01\_LIMITE\_TERMICO\_ABSOLUTO}:} si $T \ge 45.0\;\si{\celsius}$, se activa \texttt{CORTE\_DE\_EMERGENCIA}. Se cancela cualquier pulso y se anula el caudal ($0\,\%$) para evitar el colapso t\'ermico de la membrana. Prevalece sobre cualquier intervenci\'on humana.
  \item \textbf{\texttt{REGLA\_02\_LIMITE\_ECOLOGICO\_TERRITORIAL}:} si $\text{HR} < 30.0\,\%$, se activa \texttt{PAUSA\_ECOLOGICA}. Se suspende la extracci\'on ($0\,\%$) para evitar deshidratar el microclima o competir con la vegetaci\'on local durante sequ\'ias diurnas. Prevalece sobre intervenciones manuales.
  \item \textbf{\texttt{REGLA\_06\_SUPERVISION\_HUMANA}:} si existe una consigna manual de un operador, solo se ejecuta si es m\'as conservadora que la decisi\'on autom\'atica (\texttt{SUSPENDER\_EXTRACCION} lleva el caudal a $0\,\%$; \texttt{INHIBIR\_PULSO} ante alta impedancia sustituye el pulso por caudal reducido al $50\,\%$).
  \item \textbf{\texttt{REGLA\_04\_PROTECCION\_RESERVA\_CRITICA}:} si $Z > 1500.0\;\Omega$ pero $E_{\text{bat}} < 30.0\,\%$, se activa \texttt{REDUCIR\_CAUDAL\_SIN\_PULSO}. Se restringe el caudal al $50\,\%$ sin aplicar calor, evitando agotar la bater\'ia en horas de oscuridad.
  \item \textbf{\texttt{REGLA\_03\_DESPRENDIMIENTO\_SAL\_POR\_IMPEDANCIA}:} si $Z > 1500.0\;\Omega$ y $E_{\text{bat}} \ge 30.0\,\%$, se activa \texttt{APLICAR\_PULSO\_RESISTIVO}. Se dispara un pulso t\'ermico de $90\;\text{s}$ hacia la temperatura objetivo de $44.0\;\si{\celsius}$, con caudal operativo ajustado al $80\,\%$.
  \item \textbf{\texttt{REGLA\_05\_OPERACION\_NOMINAL}:} en cualquier otra condici\'on, el nodo opera en \texttt{FILTRACION\_PASIVA\_CONTINUA} al $100\,\%$ del caudal disponible, respetando el tope m\'aximo de extracci\'on ecol\'ogica ($0.33$).
\end{enumerate}

\section{Ecuaciones del modelo f\'isico continuo}
El modelo f\'isico reproduce la interacci\'on superficial y termodin\'amica mediante ecuaciones continuas compartidas entre la API y el simulador:
\begin{align}
  G_{\text{rel}}(Z) &= \min\left(1.0,\, \frac{Z_{\text{nom}}}{Z}\right), \quad Z_{\text{nom}} = 1000.0\;\Omega \label{eq:conductancia} \\
  g(T) &= \max\left(0.0,\, \min\left(1.0,\, \frac{T - 30.0}{44.0 - 30.0}\right)\right) \label{eq:factor_temp} \\
  \eta(T, Z_{\text{sal}}) &= \eta_{\text{m\'ax}} \cdot g(T) \cdot \frac{Z_{\text{ref}}}{Z_{\text{ref}} + Z_{\text{sal}}}, \quad Z_{\text{ref}} = 2500.0\;\Omega \label{eq:eficacia} \\
  Q_{\text{disp}}(\text{HR}) &= Q_{\text{m\'ax}} \cdot \left[\max\left(0.0,\, \frac{\text{HR} - 20.0}{100.0 - 20.0}\right)\right]^{1.5} \label{eq:flujo} \\
  E_{\text{pulso}}(\Delta T) &= \frac{C_{\text{matriz}} \cdot \Delta T + G_{\text{perd}} \cdot \Delta T \cdot \Delta t_{\text{pulso}}}{3600} \label{eq:energia}
\end{align}
donde $C_{\text{matriz}} = 800.0\;\text{J/K}$, $G_{\text{perd}} = 30.0\;\text{W/K}$ y $\Delta t_{\text{pulso}} = 90\;\text{s}$.

La cin\'etica de degradaci\'on de integridad de la membrana ($I$, expresada en \% con $I_0 = 100\,\%$) se calcula integrando tres t\'erminos:
\begin{equation}
  \Delta I = \Delta I_{\text{fatiga}} + \Delta I_{\text{t\'ermica}} + \Delta I_{\text{crist}}
\end{equation}
con fatiga por pulso de $\SI{0.02}{\percent}$ por evento, degradaci\'on t\'ermica de $\SI{0.05}{\percent\per(\celsius\cdot\minute)}$ para $T > 45.0\;\si{\celsius}$, y degradaci\'on por cristalizaci\'on salina continua de $\SI{0.10}{\percent\per\hour} \times (Z_{\text{sal}} / Z_{\text{nom}})$.

\section{Contrato de interfaz OpenAPI y rutas de la API}
La API publica su especificaci\'on bajo el est\'andar OpenAPI 3.1 en \texttt{docs/api/openapi.json}. La \cref{tab:rutas_api} detalla los 13 endpoints implementados en el servicio.

\begin{table}[htbp]
\centering
\small
\caption{Cat\'alogo de rutas REST de la API de Axioma-H Engine (versi\'on 0.5.0).}
\label{tab:rutas_api}
\begin{tabularx}{\textwidth}{@{}L{1.8cm}L{5.4cm}Y@{}}
\toprule
\textbf{M\'etodo} & \textbf{Ruta} & \textbf{Descripci\'on funcional} \\
\midrule
\texttt{GET} & \texttt{/} & Estado general del motor y metadatos de versi\'on. \\
\texttt{GET} & \texttt{/api/v1/health} & Verificaci\'on de salud de la API. \\
\texttt{POST} & \texttt{/api/v1/nodo/evaluar} & Evaluaci\'on de telemetr\'ia y decisi\'on puntual de reglas. \\
\texttt{POST} & \texttt{/api/v1/simulacion/ejecutar} & Simulaci\'on de 72 horas comparando reglas vs. temporizador. \\
\texttt{POST} & \texttt{/api/v1/analisis/barrido-salinidad} & Barrido en 6 escenarios con 5 semillas estoc\'asticas. \\
\texttt{POST} & \texttt{/api/v1/analisis/sensibilidad} & An\'alisis de sensibilidad global en par\'ametros clave. \\
\texttt{POST} & \texttt{/api/v1/analisis/incertidumbre} & Propagaci\'on de incertidumbre en flujo y eficacia. \\
\texttt{POST} & \texttt{/api/v1/analisis/robustez-respuesta} & Robustez del cruce de recuperaci\'on (54 combinaciones). \\
\texttt{POST} & \texttt{/api/v1/analisis/robustez-desgaste} & Robustez de reducci\'on de desgaste (18 combinaciones). \\
\texttt{POST} & \texttt{/api/v1/analisis/calibracion} & Calibraci\'on inversa del flujo contra la literatura. \\
\texttt{POST} & \texttt{/api/v1/analisis/dimensionamiento} & C\'alculo de bater\'ia de nodos seg\'un est\'andares OMS/Sphere. \\
\texttt{POST} & \texttt{/api/v1/analisis/impacto-entornos} & Evaluaci\'on en 18 entornos y comparaci\'on con ICBF. \\
\texttt{GET} & \texttt{/api/v1/sitios} & Cat\'alogo de perfiles clim\'aticos y salinos de referencia. \\
\bottomrule
\end{tabularx}
\end{table}

\section{Estructura del tablero interactivo y niveles de lectura}
El frontend (\texttt{frontend/app.py}) organiza los resultados en cinco vistas conceptuales:
\begin{itemize}[leftmargin=*]
  \item \textbf{Vista 1 (Investigaci\'on y gu\'ia):} introducci\'on pedag\'ogica, supuestos con mediciones requeridas y orientaci\'on interdisciplinaria.
  \item \textbf{Vista 2 (Estado del nodo):} consola interactiva para inyectar vectores de telemetr\'ia y visualizar la regla activada.
  \item \textbf{Vista 3 (Simulaci\'on de 72 horas):} series temporales comparativas de temperatura, humedad, impedancia, reserva y volumen acumulado.
  \item \textbf{Vista 4 (Respuesta a la pregunta):} evaluaci\'on estructurada en cuatro partes: a) reducci\'on del desgaste, b) cruce de recuperaci\'on, c) seguridad operativa y d) impacto en entornos reales.
  \item \textbf{Vista 5 (Explorador de hip\'otesis):} herramienta anal\'itica para ajustar par\'ametros propios, recalibrar el flujo y formular conjeturas.
\end{itemize}

El tablero incorpora un selector lateral de \textbf{Nivel de lectura}:
\begin{enumerate}
  \item \textbf{Comunidad:} s\'intesis en lenguaje natural sin unidades el\'ectricas, enfocada en la entrega de agua y comparaci\'on comunitaria.
  \item \textbf{Evaluador no t\'ecnico:} resumen ejecutivo enfocado en la respuesta central, decisiones de inversi\'on y seguridad de los equipos.
  \item \textbf{Investigador (por defecto):} presentaci\'on t\'ecnica exhaustiva con gr\'aficos interactivos de Altair, intervalos de confianza y acceso a los datos abiertos.
\end{enumerate}

